
\subsection{Installation and Quickstart}

HiveWatch is installed from PyPI in a single command:

\begin{lstlisting}[style=PythonStyle,language=bash,caption={Installation.}]
pip install hivewatch            # core only, no third-party deps
pip install hivewatch[all]       # with W&B, MLflow, and geolocation
\end{lstlisting}

The following listing shows a minimal integration of HiveWatch into any FL server loop, requiring no knowledge of the underlying framework:

\begin{lstlisting}[style=PythonStyle,caption={Minimal HiveWatch integration.}]
import hivewatch
from hivewatch.emitters import SSEEmitter, WandbEmitter

hivewatch.init(
    algorithm="FedAvg",
    emitters=[
        SSEEmitter(runs_dir="runs", port=7070),
        WandbEmitter(project="my-fl-project"),
    ],
)

for round_num in range(num_rounds):
    hivewatch.round_start(round_num)

    for client_id, result in client_results.items():
        hivewatch.log_client_update(
            client_id,
            round=round_num,
            local_accuracy=result["acc"],
            local_loss=result["loss"],
            num_samples=result["n"],
            bytes_sent=result["bytes"],
            train_time_sec=result["t"],
        )

    hivewatch.log_round(
        round_num,
        global_accuracy=agg_acc,
        global_loss=agg_loss,
    )

hivewatch.finish()
\end{lstlisting}

\subsection{Framework Integration}

We demonstrate framework-agnosticism by integrating the identical five-call API into three FL frameworks with different communication stacks. Full runnable examples ship with the HiveWatch source distribution.

\textit{APPFL (gRPC).} The APPFL example subclasses \texttt{ServerAgent} and intercepts each client update in \texttt{global\_update()}, calling \texttt{round\_start()} on each round transition and \texttt{log\_client\_update()} after aggregation returns. Every HiveWatch runtime call is made on the server.

\textit{Flower.} The Flower example subclasses \texttt{FedAvg}, forwarding per-client metrics carried in Flower's \texttt{FitRes} objects to \texttt{log\_client\_update()} from \texttt{aggregate\_fit()}, and logging the round summary from \texttt{aggregate\_evaluate()} once the global metrics are available. Failed clients reported by Flower are recorded through \texttt{log\_dropout()}.

\textit{NVIDIA FLARE.} The NVIDIA FLARE example subclasses \texttt{FedAvg} to inject HiveWatch calls into \texttt{send\_model()}, \texttt{\_aggregate\_one\_result()}, and \texttt{\_get\_aggregated\_result()}, and swaps the wrapped controller into the recipe's deployment map.

The bridging layer is confined to a single server-side file per framework --- approximately 75 lines for APPFL and Flower, and 115 for NVIDIA FLARE, whose controller must additionally be reconstructed and swapped into the recipe's deployment map --- and none of the three required any modification to HiveWatch itself. Clients otherwise remain framework-native: the Flower example touches no HiveWatch code on the client at all, populating only Flower's own \texttt{fit()} metrics dictionary, which the strategy then forwards. The single optional client-side addition is a call to HiveWatch's \texttt{get\_location()} helper when geographic metadata is wanted for the map dashboard; the APPFL and NVIDIA FLARE examples use it to tag each update with the client's coordinates.

\subsection{Instrumentation Overhead}

Because HiveWatch runs on the server's critical path, its overhead must be small relative to typical round durations. We measure per-call overhead over 50 rounds with 20 clients per round (1{,}050 calls total), comparing no-emitter (bookkeeping only) and SSEEmitter (synchronous JSONL append plus full \texttt{.map.json} rewrite per call) configurations. The workload is synthetic --- no model, dataset, or FL framework is involved --- so the reported times reflect instrumentation cost alone. The measurement script ships with the repository as \texttt{benchmarks/instrumentation\_overhead.py}, so the table can be regenerated on any machine. Results are shown in Table~\ref{tab:overhead}.

\begin{table}[t]
\caption{Per-call instrumentation overhead ($\mu$s), 50 rounds $\times$ 20 clients.}
\label{tab:overhead}
\centering
\begin{tabular}{lrrr}
\toprule
\textbf{Call (configuration)} & \textbf{Mean} & \textbf{Median} & \textbf{p95} \\
\midrule
\texttt{log\_client\_update} (no emitter) & 5.0 & 4.7 & 5.9 \\
\texttt{log\_round} (no emitter)           & 37.2 & 34.3 & 48.9 \\
\texttt{log\_client\_update} (SSEEmitter)  & 2822.5 & 2759.0 & 4763.8 \\
\texttt{log\_round} (SSEEmitter)           & 3005.4 & 2995.7 & 4723.7 \\
\bottomrule
\end{tabular}
\end{table}

Internal bookkeeping (schema normalization, derived-metric computation, round-state tracking) costs 5--37\,$\mu$s per call regardless of which emitters are attached. The SSEEmitter raises per-call overhead to $\sim$2.8\,ms, which is entirely attributable to the synchronous file flush after each JSONL append and the full rewrite of the map metadata snapshot — a deliberate consistency tradeoff so that crashed runs leave a complete, replayable artifact on disk. Summed over a 20-client round, the SSEEmitter adds $\sim$59\,ms of server overhead, which is negligible relative to real FL round durations (seconds to minutes). The WandbEmitter and MLflowEmitter offload persistence to external services and do not incur per-call file I/O. The SSEEmitter produced a 1.09\,MB JSONL history file and a 439\,KB map metadata snapshot for this 50-round, 20-client workload — sufficient to fully reconstruct client geography, status transitions, and round summaries via the map dashboard without re-running training.
